\section{Adaptive Fee-Burning Protocol}
\label{sec:fees}

\subsection{Transaction Fee Structure}

Lux Network uses a two-component fee model inspired by EIP-1559 \cite{roughgarden2020eip1559} but extended for multi-chain environments:
\begin{equation}
    \text{Fee}(\text{tx}) = \text{BaseFee}_t \cdot \text{GasUsed}(\text{tx}) + \text{PriorityFee}(\text{tx})
\end{equation}

The base fee adjusts per block according to utilization:
\begin{equation}\label{eq:basefee}
    \text{BaseFee}_{t+1} = \text{BaseFee}_t \cdot \left(1 + \frac{1}{8} \cdot \frac{G_t - G^*}{G^*}\right)
\end{equation}
where $G_t$ is the gas used in block $t$ and $G^* = 15{,}000{,}000$ is the target gas per block.

\subsection{Burn Rate Function}

The burn rate $\phi_t$ determines the fraction of base fees burned versus distributed to validators:
\begin{equation}\label{eq:burn}
    \phi_t = \phi_{\min} + (\phi_{\max} - \phi_{\min}) \cdot \min\left(1, \frac{G_t}{G_{\text{cap}}}\right)
\end{equation}
where $\phi_{\min} = 0.50$, $\phi_{\max} = 0.90$, and $G_{\text{cap}} = 30{,}000{,}000$. Under light load, 50\% of base fees are burned. Under maximum load, 90\% are burned. The total burn in epoch $t$ is:
\begin{equation}
    B_t = \sum_{\text{blocks } b \in t} \phi_b \cdot \text{BaseFee}_b \cdot G_b
\end{equation}

Priority fees are never burned; they go entirely to the block proposer as a tip for transaction ordering.

\begin{proposition}[Deflationary Threshold]
The network becomes net deflationary when the average block utilization exceeds $\bar{G} > G^* \cdot R_t / (\phi_{\min} \cdot \text{BaseFee}_t \cdot N_b)$, where $N_b$ is the number of blocks per epoch.
\end{proposition}

\subsection{Multi-Chain Fee Aggregation}

Each appchain $k$ contributes to the primary network burn through an appchain fee tax $\tau_k$:
\begin{equation}
    B_t^{\text{total}} = B_t^{\text{primary}} + \sum_{k \in \mathcal{K}} \tau_k \cdot F_k^t
\end{equation}
where $F_k^t$ is the total fee revenue on appchain $k$ in epoch $t$. The tax rate $\tau_k$ depends on the appchain's economic relationship with the primary network (see Section~\ref{sec:appchain}).

\subsection{Fee Burn Stabilization}

To prevent excessive deflation that could harm network usability, we introduce a floor on circulating supply:
\begin{equation}
    B_t^{\text{effective}} = \min\left(B_t^{\text{total}}, \max\left(0, S_t - S_{\text{floor}}\right)\right)
\end{equation}
where $S_{\text{floor}} = 0.5 \cdot S_0 = 250{,}000{,}000$ LUX. Once circulating supply reaches the floor, burning ceases and all fees are redistributed to validators.

